% >>> PROPAGATION (Fable 5, Aug 29 2026): the missing neighborhood placed --- the retrocausal / all-at-once programs and Adlam's
% constraint account, kinship (constraints on a block) and differentia (no un-indexed content; validity constraints on
% information states) --- as a new pre-empted objection; four bib entries added VERIFY-flagged in the local block. See R109.
% ============================================================
% DRAFTING HEADER (strip at deposit)
% Paper 7 of the corpus: the causality extraction.
% Drafted July 2026 from the Aethic long's sec:causality_without_the_triad
% (all four subsubsections) at the author's commission; self-contained
% machinery in Section 2; siblings cited via placeholder keys where the
% deposited keys are unassigned (VERIFY comments mark every such site).
% ============================================================
\documentclass[11pt]{article}
\usepackage[margin=1.05in]{geometry}
\usepackage{amsmath, amssymb, amsthm}
\usepackage{microtype}
\usepackage[hidelinks]{hyperref}
\usepackage[style=numeric, sorting=none]{biblatex}

\begin{filecontents}[overwrite]{ct_refs.bib}
% referee's read --- add it there, whereupon this local becomes a shadowed duplicate.
@comment{VERIFY at submission (entered from working knowledge, network unavailable; refs.bib canonical): wharton2015universe, price2012does, sutherland2008causally, adlam2022laws --- added Aug 29 2026 for the all-at-once placement.}
@incollection{wharton2015universe, author={Wharton, Ken}, title={The universe is not a computer}, booktitle={Questioning the Foundations of Physics: Which of Our Fundamental Assumptions Are Wrong?}, editor={Aguirre, Anthony and Foster, Brendan and Merali, Zeeya}, publisher={Springer}, pages={177--189}, year={2015}}
@article{price2012does, author={Price, Huw}, title={Does time-symmetry imply retrocausality? {How} the quantum world says ``maybe''}, journal={Studies in History and Philosophy of Modern Physics}, volume={43}, number={2}, pages={75--83}, year={2012}}
@article{sutherland2008causally, author={Sutherland, Roderick I.}, title={Causally symmetric {Bohm} model}, journal={Studies in History and Philosophy of Modern Physics}, volume={39}, number={4}, pages={782--805}, year={2008}}
@article{adlam2022laws, author={Adlam, Emily}, title={Laws of nature as constraints}, journal={Foundations of Physics}, volume={52}, pages={28}, year={2022}}
@misc{schmid2020unscrambling, author={Schmid, David and Selby, John H. and Spekkens, Robert W.}, title={Unscrambling the omelette of causation and inference: The framework of causal-inferential theories}, year={2020}, note={arXiv:2009.03297}}
% VERIFY at merge (fields from memory): The Direction of Time, posthumous.
@incollection{verma1990equivalence,
  author = {Verma, Thomas and Pearl, Judea},
  title = {Equivalence and synthesis of causal models},
  booktitle = {Proceedings of the Sixth Conference on Uncertainty in Artificial Intelligence},
  pages = {255--268}, year = {1990}
}
@book{spirtes2000causation,
  author = {Spirtes, Peter and Glymour, Clark and Scheines, Richard},
  title = {Causation, Prediction, and Search},
  edition = {2}, publisher = {MIT Press}, address = {Cambridge, MA}, year = {2000}
}
@book{pearl2009causality,
  author = {Pearl, Judea},
  title = {Causality: Models, Reasoning, and Inference},
  edition = {2}, publisher = {Cambridge University Press}, address = {Cambridge}, year = {2009}
}
@misc{AethusRealist, author={Benander, Ajax}, title={Determinacy as a Two-Place Relation: The Category Mistake in Quantum Interpretation, and the Dissolution of the Paradox Set}, year={2026}, howpublished={Preprint}, url={https://philpapers.org/rec/BENTRO-72}}
@misc{AethusSemantics, author={Benander, Ajax}, title={Counterfactual Evaluation by Perturbation of Information States: Deleting the World-Similarity Ordering, Not Refining It}, year={2026}, howpublished={Preprint}, url={https://philpapers.org/rec/BENCEB-2}}

@article{galles1998axiomatic, author={Galles, David and Pearl, Judea}, title={An Axiomatic Characterization of Causal Counterfactuals}, journal={Foundations of Science}, volume={3}, number={1}, pages={151--182}, year={1998}}

@book{hoferszabo2013principle,
  author = {Hofer-Szab\'o, G\'abor and R\'edei, Mikl\'os and Szab\'o, L\'aszl\'o E.},
  title = {The Principle of the Common Cause},
  publisher = {Cambridge University Press}, address = {Cambridge}, year = {2013}
}
@book{reichenbach1956direction,
  author = {Reichenbach, Hans},
  title = {The Direction of Time},
  publisher = {University of California Press}, year = {1956}}
% VERIFY at merge (fields from memory; PRX 7, 031021 believed correct).
@article{allen2017quantum,
  author = {Allen, John-Mark A. and Barrett, Jonathan and Horsman, Dominic C. and Lee, Ciar\'an M. and Spekkens, Robert W.},
  title = {Quantum Common Causes and Quantum Causal Models},
  journal = {Physical Review X}, volume = {7}, year = {2017}}
% VERIFY at merge (fields from memory): the tickle-defense home.
@book{eells1982rational,
  author = {Eells, Ellery},
  title = {Rational Decision and Causality},
  publisher = {Cambridge University Press}, year = {1982}}
% VERIFY at merge (fields from memory): the medical-Newcomb discussions.
@book{skyrms1980causal,
  author = {Skyrms, Brian},
  title = {Causal Necessity: A Pragmatic Investigation of the Necessity of Laws},
  publisher = {Yale University Press}, year = {1980}}
% via the IfFileExists line, eliminating the same-key collision risk at merge.
% AethusBrief's title below is verbatim from the Brief's own title page.
@misc{AethusBrief, author={Benander, Ajax}, title={Aethic Reasoning in Brief: From a Counterfactual Semantics to the Quantum Measurement Problem}, year={2026}, howpublished={Preprint}, url={https://philpapers.org/rec/BENARI-3}, eprint={https://philsci-archive.pitt.edu/id/eprint/30439}, eprinttype={Alternative URL}}
% VERIFY at merge: title verbatim from the deposited boundary extraction; fill URL.

% VERIFY at merge: fill the deposit URL once the decision extraction is deposited.
@misc{AethusDecision, author={Benander, Ajax}, title={Newcomb's Problem Without a Fixed Past: A Decision Rule From Observer-Indexed Determinacy}, year={2026}, howpublished={Preprint}, url={https://philpapers.org/rec/BENNPW}}
@misc{Aethus, author={Benander, Ajax}, title={Aethic Reasoning}, year={2026}, note={Companion manuscript; deposit locator to be added}, url={https://philpapers.org/rec/BENARA-5}}

@misc{AethusQuant,
  author       = {Benander, Ajax},
  title        = {Quantum Mechanics from Records: Measurement Structure on a Finite Path-Kernel Domain},
  year         = {2026},
  howpublished = {Preprint},
  url          = {https://philpapers.org/rec/BENTQB},
  note         = {The physics-face axiomatization; first posted July 17, 2026.}
}

@article{fine1982hidden,
  author = {Fine, Arthur},
  title = {Hidden variables, joint probability, and the {Bell} inequalities},
  journal = {Physical Review Letters},
  volume = {48}, number = {5}, pages = {291--295}, year = {1982}
}

@incollection{jaynes1990probability, author={Jaynes, E. T.}, title={Probability in Quantum Theory}, booktitle={Complexity, Entropy, and the Physics of Information}, editor={Zurek, W. H.}, publisher={Addison-Wesley}, year={1990}}

@incollection{nozick1969newcomb, author={Nozick, Robert}, title={Newcomb's Problem and Two Principles of Choice}, booktitle={Essays in Honor of Carl G. Hempel}, editor={Rescher, Nicholas}, publisher={Reidel}, year={1969}}

@article{pitowsky1994george,
  author = {Pitowsky, Itamar},
  title = {George {Boole}'s `conditions of possible experience' and the quantum puzzle},
  journal = {British Journal for the Philosophy of Science},
  volume = {45}, number = {1}, pages = {95--125}, year = {1994}
}

@misc{spekkens2024escapedsapiens,
  author       = {Spekkens, Robert W.},
  title        = {What is Reality Made Of? {Unscrambling} the Quantum Omelet},
  howpublished = {Interview by Shane Farnsworth, The Escaped Sapiens Podcast, Episode 74},
  year         = {2024},
  month        = nov,
  note         = {Video available at \url{https://youtu.be/J2ZIRkfrFlI}},
  url          = {https://youtu.be/J2ZIRkfrFlI}
}

@article{wood_spekkens_2015, author={Wood, Christopher J. and Spekkens, Robert W.}, title={The lesson of causal discovery algorithms for quantum correlations: causal explanations of Bell-inequality violations require fine-tuning}, journal={New Journal of Physics}, volume={17}, pages={033002}, year={2015}}

@book{lauritzen1996graphical,
  author = {Lauritzen, Steffen L.},
  title = {Graphical Models},
  publisher = {Clarendon Press}, address = {Oxford}, year = {1996}
}
@book{hume1748enquiry,
  author = {Hume, David}, title = {An Enquiry Concerning Human Understanding},
  publisher = {A. Millar}, address = {London}, year = {1748}
}
@article{mackie1965causes,
  author = {Mackie, J. L.}, title = {Causes and Conditions},
  journal = {American Philosophical Quarterly}, volume = {2}, number = {4}, pages = {245--264}, year = {1965}
}
@book{mackie1974cement,
  author = {Mackie, J. L.}, title = {The Cement of the Universe: A Study of Causation},
  publisher = {Clarendon Press}, address = {Oxford}, year = {1974}
}
@book{suppes1970probabilistic,
  author = {Suppes, Patrick}, title = {A Probabilistic Theory of Causality},
  publisher = {North-Holland}, address = {Amsterdam}, year = {1970}
}
@book{eells1991probabilistic,
  author = {Eells, Ellery}, title = {Probabilistic Causality},
  publisher = {Cambridge University Press}, address = {Cambridge}, year = {1991}
}
@article{cartwright1979causal,
  author = {Cartwright, Nancy}, title = {Causal Laws and Effective Strategies},
  journal = {No{\^u}s}, volume = {13}, number = {4}, pages = {419--437}, year = {1979}
}
@article{simpson1951interpretation,
  author = {Simpson, E. H.}, title = {The Interpretation of Interaction in Contingency Tables},
  journal = {Journal of the Royal Statistical Society, Series B}, volume = {13}, number = {2}, pages = {238--241}, year = {1951}
}
@book{salmon1984scientific,
  author = {Salmon, Wesley C.}, title = {Scientific Explanation and the Causal Structure of the World},
  publisher = {Princeton University Press}, address = {Princeton}, year = {1984}
}
@book{dowe2000physical,
  author = {Dowe, Phil}, title = {Physical Causation},
  publisher = {Cambridge University Press}, address = {Cambridge}, year = {2000}
}
@article{lewis1973causation,
  author = {Lewis, David}, title = {Causation},
  journal = {Journal of Philosophy}, volume = {70}, number = {17}, pages = {556--567}, year = {1973}
}
@book{halpern2016actual,
  author = {Halpern, Joseph Y.}, title = {Actual Causality},
  publisher = {MIT Press}, address = {Cambridge, MA}, year = {2016}
}
@book{woodward2003making,
  author = {Woodward, James}, title = {Making Things Happen: A Theory of Causal Explanation},
  publisher = {Oxford University Press}, address = {New York}, year = {2003}
}

@misc{GoldenScience2022,
  author       = {Benander, Ajax},
  title        = {Golden Science},
  year         = {2022},
  month        = {November},
  howpublished = {Self-hosted manuscript},
  url          = {https://aethus.org/aethic-code/files/GoldenScience_Nov2022.pdf},
  note         = {Title page August 2022, the date of the first version of the source; this version created November 2022. Not deposited to a
                  preprint server; cited as a dated developmental record rather than
                  as a source of results.}
}
\end{filecontents}
\IfFileExists{refs.bib}{\addbibresource{refs.bib}}{} % corpus bibliography where present; refs.bib wins collisions
\addbibresource{ct_refs.bib} % local standalone coverage

\newtheorem{definition}{Definition}
\newtheorem{theorem}{Theorem}
\newtheorem{proposition}{Proposition}
\newtheorem{corollary}{Corollary}
\newtheorem{principle}{Principle}
\newtheorem{example}{Example}

\title{Causality from Indexed Information States:\\ \large Screening, Intervention, and the Bell Obstruction}
\author{Ajax Benander\thanks{Email: \href{mailto:akb9408@rit.edu}{akb9408@rit.edu}, Website: \url{https://aethus.org}.}}
\date{First drafted July 2026\\ {\small This version: October 2026 --- working draft}}

\begin{document}
\maketitle

\begin{abstract}
This paper develops a representation of statistical dependence on observer-indexed partial information states and determines what it recovers of causal reasoning. Three levels of result are distinguished. The first is observational. For factorized states in the regime of mutually exclusive alternatives, a correlation forces connectivity through graded couplings, separating families screen, and every screener carries the whole covariance; this is the statistical screening structure associated with Reichenbach's common-cause principle, located in graphical sufficient conditions, factorization equivalences and covariance identities. The second is conditional on supplied structure. Given directed mechanism histories with independent mechanism weights, settling an attribute is a surgery and the interventional calculus is inherited; given functional mechanism entries, counterfactual evaluation coincides with the functional-model one. A representation theorem shows that hard invalidity together with independent auxiliary weights, constrained and renormalized, represents every exclusive-regime state; the construction is distinct from a directed mechanism history and supplies no direction. The third is a boundary. In the regime of agreeing alternatives, Bell-violating statistics admit no global noncontextual, measurement-independent valuation extending the observed table, while the observational results stand on that table. The framework's record discipline, stated as a premise, supplies its proposed explanation of why no admissible state carries such a mediator; it adds a restriction on admissible states, not a further statistical exclusion. The construction motivates an account of causal dependence in terms of indexed coupling. Causal orientation remains additional structure that the framework supplies by convention and does not derive.
\end{abstract}

\section{Introduction}
\label{sec:intro}

Reichenbach's Common Cause Principle \cite{reichenbach1956direction} makes two claims at once: a \emph{demand}, that every improbable correlation be grounded, and a \emph{schema} for the grounding, the disjunction of three producer stories with the common-cause horn certified by screening-off, conditional independence of the correlates given the cause. The Bell correlations strain the pair, and Wood and Spekkens \cite{wood_spekkens_2015} converted the strain into a theorem: any classical causal model recovering the quantum statistics must fine-tune its parameters to hide the causal pathways it posits. The constructive wing of the same program generalizes the causal models themselves, retaining directed causal structure and paying for Bell in quantum conditional states \cite{allen2017quantum}; the diagnostic wing asks what causal explanation should be once the triad's exhaustiveness is in doubt. This paper belongs to the second wing. Its primitive is not a producer relation between events but a \emph{jointness} relation among attribute-values relative to an information state: graded couplings recording which combinations of alternatives stand together, and at what weight.

\noindent\textit{The motivating case.} A driver is stopped at a red light; two minutes later a sinkhole opens just ahead, where the car would have been on a green. Relative to the driver's later state, which settles the sinkhole and its underground history, the counterfactual holds: had the light been green, the car would have met the hole. Relative to the state at the light it does not, because there the sinkhole and its history are open, and that state's continuations include ones in which the car crosses the same ground with no hole beneath it, whichever the light. With $G$ the green light and $H$ meeting the hole, $A_{\mathrm{later}} \models G \mathrel{\Box\!\!\rightarrow} H$ while $A_{\mathrm{earlier}} \not\models G \mathrel{\Box\!\!\rightarrow} H$. The later verdict does not overturn the earlier one: the earlier openness remains true relative to the earlier state, and the erosion's having preceded the light does not supply, at that index, the determination the earlier state lacks, temporal location and determinacy at an index being separate relations. On the proposed interpretation, the earlier state is complete as a description of reality relative to that index; its openness is not an incomplete description of an independently determinate history. The semantics companion establishes the indexed difference between the two verdicts \cite{AethusSemantics}, and the determinacy companion defends the ontological reading \cite{AethusRealist}. A later state's categorical dependence verdict is therefore not exported backward into an earlier state merely because the relevant events lie in its past, and that export is the step the causal-chain picture takes when it reads every dependence off one completely specified history. Dependence is also selective: a script that buys breakfast according to whether the light was red installs a dependence of breakfast on the light, relative to the state in which the script was written and not relative to its siblings in which it was not, and the cat on the landing, stipulated to stand in no dependency on the light, is untouched by it. A suitable state can be constructed for each dependency separately (Proposition~\ref{prop:tie}), but that does not establish a single state containing all of them, $\forall Y\,\exists A_Y\,\mathrm{Dep}_{A_Y}(Y, X) \not\Rightarrow \exists A\,\forall Y\,\mathrm{Dep}_A(Y, X)$, and the familiar argument that a changed light changes everything after it slides between such states as though it traced one fixed structure \cite{Aethus}. Here and below $\mathrm{Dep}_A(Y, X)$ denotes a dependence carried by $A$'s own entries, a constraint of $A$ touching $X$ and $Y$ that restricts or weights their joint combinations; Section~\ref{sec:framework} separates this structural relation from the statistical, interventional and reopening relations the paper also uses, and the tie proposition establishes the structural one without identifying any interventional direction.

Three propositions are kept apart below. \emph{Indexing}: determinate content is relative to an information state. \emph{Selective dependence}: a perturbation need not alter every other attribute; this has classical counterparts, and it is no argument against deterministic dynamics, since a deterministic law can leave one variable untouched by a change to another. \emph{Tense-free openness}: a past event can remain constructively open relative to an index that does not settle it. The first and third carry the proposal. Its target is the causal-chain picture, dependence read off a complete determinate history taken as the ontology of what a state leaves open, and not deterministic dynamics or the practice of causal modeling as such.

\noindent\textit{Why a global history cannot be assembled by transfer.} Suppose $\mathrm{Dep}_A(Y, X)$ and $\mathrm{Dep}_B(Z, Y)$. Neither fact alone authorizes a composed dependence of $Z$ on $X$. Before the two compose, a common admissible index must be established, both dependencies must be shown to survive transport to it, and the notion of dependence in use needs a composition rule, dependence not being automatically transitive. The causal-chain picture builds its determinate history by composing dependencies this way without those steps, each link taken from whatever index supplied it and ordered by the temporal order of the events. The framework blocks that unlicensed construction. It does not prohibit a common mechanism history wherever the compatibility and composition conditions actually hold.

\noindent\textit{Two counterfactual operations.} The motivating verdicts are reopening verdicts, the semantics companion's operation \cite{AethusSemantics}: reopening retracts the antecedent's entry and the content grounded on it before restating the antecedent. The counterfactual result of Section~\ref{sec:rungs} concerns a different operation, functional-model settling, which keeps the inferred background distribution and performs surgery on the mechanism. The two are not interchangeable in general; this paper proves its correspondence result for settling alone (Proposition~\ref{prop:cf}) and claims no general agreement between settling and reopening.

\noindent\textit{Results, in three groups.} (1) For factorized observational states in the regime of mutually exclusive alternatives: connectivity through graded couplings, separation, screening as pairwise composition, and the covariance identity, the statistical consequences associated with Reichenbach's principle (Section~\ref{sec:theorem}). A separator carrying the covariance is not thereby shown to be a prior common cause, and the fork inequalities are the binary case of the covariance identity. (2) Given a directed mechanism history, supplied with the state and not derived from it: settling is Pearl's truncated factorization, so the interventional calculus is inherited, and counterfactual evaluation coincides with the functional-model one where functional structure over background attributes is supplied (Section~\ref{sec:rungs}). A representation theorem writes every exclusive-regime state as hard constraints over independent auxiliary weights, constrained and renormalized, a construction distinct from a mechanism history and carrying no direction (Section~\ref{sec:rerendition}). (3) In the regime of agreeing alternatives, Bell-violating statistics obstruct a specified global valuation extension, noncontextual and measurement-independent, while the observational results stand on the observed table (Section~\ref{sec:boundary}); the record discipline, imported as a premise, supplies the framework's explanation of why no admissible state carries such a mediator, a restriction on admissible states and not a further statistical exclusion. Beneath the three stands one principle, stated in Section~\ref{sec:principle}: correlation is represented by joint coupling, joint invalidity its structural layer and weights grading what that layer admits, with no direction and no mechanism of its own.

\noindent\textit{Contribution and limits.} The screening steps of (1) are standard graphical-model results, and no new probability identity is claimed for them. The paper accepts directed input structure as supplied and proves a correspondence with it; what it declines is the reading of that structure as primitive production, and whether that reading can be declined is a philosophical question the paper argues rather than settles. What the formulation adds is a different ontology for structures that conventional modeling already contains: an indexed state describes reality relative to its index completely, its blanks are not missing entries from an additional determinate history, and a later state's categorical verdicts are not exported to an earlier one. Section~\ref{sec:rerendition} exhibits one consequence of that ontology which a probability table does not carry, two states with one table and different hard structure, and follows one state through observation, direction, intervention and counterfactual, each step adding one structure the previous step lacked. Causal orientation is supplied by convention and not derived (Section~\ref{sec:tensefree}); couplings alone fix no interventional law (Proposition~\ref{prop:parity}); conditional tables fix no counterfactual (Proposition~\ref{prop:cf}). The framework the primitive comes from is developed elsewhere \cite{AethusBrief, Aethus}; Section~\ref{sec:framework} states the machinery this paper needs, and the quantum classification \cite{AethusQuant} and the decision extraction \cite{AethusDecision} are engaged where their results are load-bearing and not re-argued. The enumerations reported beside the theorems are reproduced by the supplementary scripts described in Section~\ref{sec:conclusion}.

\section{The machinery, minimally}
\label{sec:framework}

An \emph{information state} assigns values to some attributes and leaves others open.\footnote{The fuller development this paper draws on \cite{Aethus} gives the states a proper name, \emph{Aethus}, plural \emph{Aethae}; the present paper needs only the structure and keeps the generic vocabulary.} An attribute carries an \emph{answer space} of standing alternatives; an attribute with a value is \emph{stated}, one without is \emph{blank}. Updating is multiplicative: stating an attribute multiplies its factor into the state, prunes what the new content excludes, and renormalizes --- so a state's constraint content is a product of entry-factors over the attribute families each entry touches.

Two precisions carry the paper. First, entries carry two layers. A structural layer of \emph{joint invalidity} determines which combinations of alternatives across the families an entry touches can stand together at all; a weight layer grades the combinations that can. Hard invalidity determines admissibility, and weights quantify alternatives within the admissible structure: a probabilistic representation may assign zero weight to forbidden combinations, but zero weight alone does not identify structural invalidity, and the quantum companion's admissible cells of zero kernel weight are the standing example \cite{AethusQuant}. Second, the alternatives of a blank attribute can stand in either of two regimes, and the framework's canonical names for them are used here because the theorem they enter is a causal one in which coherence is not the subject. In \emph{disagreeing superposition}---the alternatives are mutually exclusive and jointly exhaustive, each stating the others' negations in its interior, exactly one to be realized; conditioning on such an attribute lands the state inside a single \emph{pocket}, and a mixture-level quantity over the whole attribute is an averaging projection over its pockets. In \emph{agreeing superposition}---the alternatives stand together with no mutual negation stated, the regime quantum superposition occupies on the fuller development's analysis \cite{Aethus}, and there are no mutually exclusive pockets to condition into. The distinction is logical, a matter of what the alternatives' interiors state, and is defined before any weight or phase; it is not the coherent/decoherent distinction of the physical realization, and the two are not identified here or elsewhere in the corpus. Disagreeing superposition is the classical regime, and Sections~\ref{sec:theorem}--\ref{sec:instances} work entirely inside it; agreeing superposition enters only at Section~\ref{sec:boundary}, as the boundary.

One more piece of vocabulary: a \emph{mixture-level correlation} between attribute families $A$ and $B$ is a dependence visible in the state's weights before conditioning. The question of the paper is what grounds it, and the answer consults no arrow.

\begin{definition}[The model]\label{def:model}
A finite set of attributes, each with a finite answer space; a finite set of \emph{entries} $e$, each with a scope $F_e$ of families, a hard part $\mathrm{Adm}_e$, the combinations of its families' alternatives it admits, and a weight $W_e \ge 0$ on those combinations; and the state's law
\[
P_A(v) \;=\; \frac{\mathbf{1}[v \in \mathrm{Adm}_A]}{Z_A}\, \prod_e W_e(v_{F_e}), \qquad \mathrm{Adm}_A = \{\,v : v_{F_e} \in \mathrm{Adm}_e \text{ for every } e\,\},
\]
closed under the framework's invalidation rules, with $Z_A > 0$. The separation structure of Section~\ref{sec:theorem} is drawn over all entries, hard ones included. This matters because pruning itself induces dependence: two independent fair bits restricted to $X = Y$ are dependent, and their unary weight factors are disconnected unless the equality constraint is an entry of the graph. The factorization theorems below therefore concern the complete final law, every hard restriction represented.
\end{definition}

\noindent Four relations go under the word \emph{dependence}, and they are distinguished once.
\begin{center}\small
\begin{tabular}{p{0.33\textwidth} p{0.57\textwidth}}
\hline
\textbf{Relation} & \textbf{What establishes it} \\
\hline
Statistical dependence & a nonfactorizing law $P_A$ \\
Structural dependence, $\mathrm{Dep}_A$ & an entry of $A$ restricting or weighting the joint combinations of its relata \\
Interventional dependence & a response under a specified settling operation (Section~\ref{sec:rungs}) \\
Reopening dependence & a verdict of the companion's modal operator \cite{AethusSemantics} \\
\hline
\end{tabular}
\end{center}
\noindent The tie proposition below establishes structural dependence; it identifies no interventional direction, which the mechanism histories of Section~\ref{sec:rungs} supply.

\begin{proposition}[Ties can be constructed]\label{prop:tie}
Work in a finite constraint model in which every terminal descendant is a total assignment and every state has a terminal descendant, the admissible total assignments extending a state's stated content are declared as its terminal refinements and are structurally visible, and the base contradictions are the assignments violating an exclusion. In such a model a state is valid exactly when some total assignment extending its stated content violates none of its exclusions, by the direct form of the third postulate \cite{Aethus}; outside it the equivalence can fail, since a consistent total assignment need not be an accessible, visible proper descendant able to supply a validity floor. Let $A$ be a valid state with binary attributes $X$ and $Y$, and let $A_{\mathrm{tie}}$ be $A$ with one entry added that excludes the combinations $(X = 0, Y = 1)$ and $(X = 1, Y = 0)$. Then $A_{\mathrm{tie}}$ is valid if and only if some valid total assignment of $A$ has $X = Y$; in every valid total assignment of $A_{\mathrm{tie}}$, $X = Y$, the mismatches excluded outright rather than weighted toward zero; and if $A$ has valid total assignments with $X = Y = 0$ and with $X = Y = 1$, the tie leaves $X$ open in $A_{\mathrm{tie}}$, both values standing.
\end{proposition}
\begin{proof}
The valid total assignments of $A_{\mathrm{tie}}$ are exactly those of $A$ with $X = Y$, and each clause follows.
\end{proof}

\noindent In such a model a tie between two attributes can be installed wherever their matching combinations are jointly admissible, and only there: attributes already fixed incompatibly by the retained constraints cannot be tied while validity is kept. This is the sense in which the joint-invalidity structure carries more than correlation. A correlation theory sees only a dependence in the weights, as ice-cream sales and shark attacks correlate whether or not either responds to the other; hardness separates a structural exclusion from a small weight, and the decomposition of the structure into entries does the rest: with a shared condition $T$, the entries $I = T$ and $S = T$ exclude the mismatches between ice-cream sales $I$ and shark attacks $S$ as outright as a direct entry on $(I, S)$ would, and the two presentations differ in which entries carry the exclusion, the shared-condition one supplying $T$ as a graphical separator and the direct one supplying none. Neither observation establishes causal direction, which the mechanism histories of Section~\ref{sec:rungs} supply as additional structure; whether a mechanism history can be extracted from the richer constraint structure is open, needing an extraction rule and its justification.

\section{Constraint representation and mechanism structure}\label{sec:rerendition}

This section names three structures a state carries and keeps them apart, because they support different conclusions; it then follows one state through the four steps those structures license, and exhibits what the hard structure carries that a probability table does not.

\begin{definition}[Three structures]\label{def:structures}
For a state presented as a product of entries:
\begin{enumerate}
\item[(G1)] its \emph{factor graph} records which families share a graded entry;
\item[(G2)] its \emph{exclusion structure} records which combinations of alternatives are forbidden, the hard parts of its entries closed under the framework's invalidation rules; it is symmetric and carries no order;
\item[(G3)] a \emph{mechanism history}, when the state has one, records which mechanisms take which families as inputs: normalized, directed entries of the kind Theorem~\ref{thm:surgery} assumes.
\end{enumerate}
\end{definition}

\noindent They support different conclusions. Correlation is weight, a dependence in the soft parts of the entries; but the connection a correlation forces runs through (G1), not (G2). The demand of Theorem~\ref{thm:reichenbach} establishes a chain of entries, soft factors included, and it cannot be strengthened to a chain of zero sets: the single entry $W(0,0) = W(1,1) = 0.4$, $W(0,1) = W(1,0) = 0.1$ is correlated, strictly positive, and has an empty hard part. Direction, in turn, is (G3), not the order of inquiry. Commuting observational updates can be made in any order without changing the state, and an earlier event can be recorded later than a more recent one, so the order in which entries were acquired carries no causal information; the settling theorem is defensible exactly because it assumes input sets, and it establishes a correspondence with causal Bayesian networks rather than deriving their direction from an update history.

\subsection{One state through the four steps}\label{sec:worked}

Take three binary attributes and one factorized observational state, $W(c, a, b) = w(c)\,w(a \mid c)\,w(b \mid c)$, with $w(c) = \tfrac12$, $w(a \mid c) = \tfrac45$, $w(a \mid \bar{c}) = \tfrac15$, $w(b \mid c) = \tfrac{7}{10}$ and $w(b \mid \bar{c}) = \tfrac{3}{10}$.

\textit{Observation.} $A$ and $B$ are correlated, $\mathrm{Cov}(A, B) = P(c)\,P(\bar{c})\,\Delta_A \Delta_B = \tfrac14 \cdot \tfrac35 \cdot \tfrac25 = \tfrac{3}{50}$, and $C$ separates and screens them (Theorem~\ref{thm:reichenbach}). Nothing yet says which way anything runs.

\textit{Direction.} Two mechanism histories carry exactly these weights: the fork, with entries $e_C(c)$, $e_A(a \mid c)$ and $e_B(b \mid c)$; and the chain $A \to C \to B$, with $e_A(a) = \tfrac12$, $e_C(c \mid a) = \tfrac45$, $e_C(c \mid \bar{a}) = \tfrac15$ and the same $e_B$. Choosing one is the added directional structure.

\textit{Intervention.} Settling $A = 1$ by the truncated factorization (Theorem~\ref{thm:settling}) gives $P(b \mid do(a)) = \tfrac12$ in the fork and $\tfrac{31}{50}$ in the chain, the latter equal to $P(b \mid a)$: one observational state, two interventional laws, as Proposition~\ref{prop:parity} requires.

\textit{Counterfactual.} In the fork, give $B$ a deterministic entry over background attributes uniform on ten values, in two ways: one shared threshold, $B = \mathbf{1}[u < 7]$ when $C = 1$ and $\mathbf{1}[u < 3]$ when $C = 0$, with $u$ uniform on $\{0, \ldots, 9\}$; or two background attributes $u_1$ and $u_0$, independent and each uniform on $\{0, \ldots, 9\}$, with $B = \mathbf{1}[u_1 < 7]$ when $C = 1$ and $\mathbf{1}[u_0 < 3]$ when $C = 0$. Both reproduce $e_B$, and both give $P(b \mid do(\bar{c})) = \tfrac{3}{10}$. For a unit with $C = 1$ and $B = 1$, the counterfactual weight of $B = 1$ had $C$ been $0$ is $\tfrac37$ under the first and $\tfrac{3}{10}$ under the second (Proposition~\ref{prop:cf}).

Each step added one structure the previous step lacked, and each addition changed what could be computed: observational weights, mechanism orientation and cross-counterfactual structure are three different inputs.

\begin{theorem}[Hard-constraint representation of factorized weights]\label{thm:snf}
Let a state in disagreeing superposition on finite attribute spaces be presented as a product of entries with nonzero normalization. Then there is an enlarged attribute space, the original families together with one auxiliary attribute $M_k$ per entry, on which every entry is a hard constraint and the state's weights are the pruning, by the compatibility constraints, of a product of independent weights $\mu_k$ on the $M_k$, renormalized. Both layers of the original state are recovered: its admissible combinations, those of weight zero included, correspond one to one with the admissible pockets of the enlarged space, and marginalizing out the $M_k$ returns its weights. Preservation of proper-child descent is a further question, not claimed here.
\end{theorem}

\begin{proof}
Let the $k$th entry weight the combinations $c$ of its families' alternatives by $W_k(c)$. Give $M_k$ the answer space of the entry's structurally admissible combinations, those its exclusions do not forbid, whatever their weight, zero included, and declare $M_k = c$ jointly invalid with every alternative of those families that $c$ omits, so the entry becomes the constraint ``$M_k = c$ fixes its families at $c$'' with weights $\mu_k(c) = W_k(c)$ on $M_k$. Admissibility is preserved: an original combination is admissible exactly when each entry's restriction $c_k$ of it is, that is, exactly when $M_k = c_k$ is available for every $k$, and the admissible pockets of the enlarged space are exactly these, each original admissible combination determining the $M_k$ uniquely. Weights are recovered: the state's constraint content being a product of entry-factors, its weight on an admissible pocket of the enlarged space is $\prod_k \mu_k(c_k) = \prod_k W_k(c_k)$, zero where some factor is zero, renormalized by the same nonzero normalization, so marginalizing over the $M_k$ returns $W$.
\end{proof}

\noindent What the representation shows is that hard invalidity together with independent auxiliary weights, subsequently constrained and renormalized, suffices to represent every exclusive-regime state. What it does not show is causal structure, and the construction is distinct from a mechanism history, which consists of directed, normalized mechanisms with specified input sets. The representation is not unique; the retained auxiliary attributes are in general dependent, since the constraints $X = M_1$ and $X = M_2$ on two independent fair bits retain only $M_1 = M_2$; some auxiliary assignments admit no state at all; and nothing in it is directed, so it carries no intervention semantics. A recursive functional model assigns every endogenous value from every background assignment, and the independence of its background variables is a property of the model, not of a representation. The representation therefore sits beside the results on factorized observational states, as a way of writing their weights, and it does not carry the results on mechanism histories, which supply the structure it lacks.

\begin{example}[One table, two hard structures]\label{ex:twostates}
Let $X$ and $Y$ be fair bits, and compare two states. In $A_{\mathrm{soft}}$ a single entry on $(X, Y)$ admits all four combinations and weights the mismatches $(0,1)$ and $(1,0)$ at zero; in $A_{\mathrm{hard}}$ the tie entry of Proposition~\ref{prop:tie} excludes the mismatches outright. The two states have one probability table, $P(0,0) = P(1,1) = \tfrac12$ and zero elsewhere, and every statistical quantity agrees. They differ in what they admit. The probe state that adds $X = 0, Y = 1$ to the stated content is admissible in $A_{\mathrm{soft}}$, a valid state of zero weight, and structurally invalid in $A_{\mathrm{hard}}$, a base contradiction; the representation of Theorem~\ref{thm:snf} records the difference in the auxiliary attribute's answer space, which contains the mismatches at $\mu = 0$ for the first state and omits them for the second. The difference has consequences. Add to each state an entry that excludes the matches $(0,0)$ and $(1,1)$. The admissibility set of $A_{\mathrm{soft}}$ remains nonempty, the mismatches surviving as structurally admissible total assignments; that of $A_{\mathrm{hard}}$ becomes empty, a base contradiction. Two states with one table thus absorb the same further constraint differently, which is the sense in which the table does not exhaust the structure. The comparison is structural and stops there. The surviving assignments of $A_{\mathrm{soft}}$ carry zero weight, so the weighted law of Definition~\ref{def:model} has normalization $Z = \sum_{x \neq y} W(x, y) = 0$ and returns no conditional probability: a further multiplicative entry cannot raise a zero weight, and replacing the original weights would be a separate operation, not the constraint update. Nor is structural admissibility validity in the weighted sense. A component carried at negligible relative weight is no proper child, so the zero-weight probe, admissible in $A_{\mathrm{soft}}$ and intrinsically valid in the abstract constraint model, supplies no validity floor for this weighted parent; $A_{\mathrm{soft}}$'s validity rests on its positive-weight refinements, and whether it has any after the exclusion is a question the example does not answer.
\end{example}

\textit{The readings, and what is derived.} Given a mechanism history, the classical readings are read off its directed structure as in a causal Bayesian network: a direct arrow; a mediating chain $A \to C \to B$, in which $C$ separates the endpoints in (G1) and screens them; a common cause, prior to both; a common effect, a mechanism taking $A$ and $B$ as inputs, which does not separate them in (G1) and conditioning on whose stated output can induce dependence; and compound separators in mixed positions. None of the readings is derived from (G1) or (G2). What is derived is this: in the exclusive regime the statistical screening structure associated with Reichenbach's principle, Theorem~\ref{thm:reichenbach}; given a mechanism history, the interventional and counterfactual correspondences, Section~\ref{sec:rungs}; and, where the two-wing statistics violate the CHSH bound, the exclusion of a global exclusive representation satisfying (D1)--(D3), Theorem~\ref{thm:reductio}, which is the exact form of the obstruction and not the unavailability of every exclusive representation of every agreeing state. Causal common ancestry is not derived, and the direction of (G3) is consumed.

\begin{proposition}[Support obstruction]\label{prop:support}
If $P(A = a) > 0$, $P(B = b) > 0$ and $P(A = a, B = b) = 0$, then $A$ and $B$ are dependent.
\end{proposition}

\begin{proof}
Independence would give $P(A = a, B = b) = P(A = a)\,P(B = b) > 0$.
\end{proof}

\noindent It is a zero in the marginal support that forces dependence, not a hard coupling touching both families: $B = A \oplus U$ with $A$ and $U$ independent fair bits is a deterministic mechanism touching $A$ and $B$, changing $B$ when $A$ changes with $U$ held, and $A \perp B$. Causal structure and observed dependence are linked only under faithfulness-type assumptions, which this paper does not add; and causation without correlation is not regime-bound, the agreeing regime merely adding cancellation at the amplitude level, the dark fringe of the double slit, as a further route to it.

\subsection{The tense-free substrate, and the conventions of precedence}\label{sec:tensefree}

Nothing in (G2) is temporal: a joint invalidity says that two determinations do not co-obtain, and it has no direction. Direction is (G3), a structure the framework consumes, and the framework reads blankness tense-free, an unrecorded past event being as open as an unhappened future one. The classical conviction that causes precede their effects in proper time adds, over one prerequisite, four commitments to that substrate, itemized so that each can be declined separately. (K0) The prerequisite: a specified functional model, mechanism entries with their input sets, which Proposition~\ref{prop:cf} shows counterfactual evaluation to need and which the observational state does not fix; not uniqueness among the models compatible with the observations, a directed structure admitting several topological orderings. (K1) \emph{Temporal alignment}, that the direction of the mechanism history agrees with proper time, an assumption and not a consequence: that records follow their events does not yield it, an earlier event being recordable later than a later one, and on the fuller development's account the record asymmetry is a statistical, observer-relative accessibility gradient, steep for entropy-reliant attributes and shallow for low-entropy ones, which the convention reads as a categorical alignment \cite{Aethus}; where the gradient is shallow, an eclipse predicted before it happens or a predictor's record of an act not yet taken, the mechanism history has the orbital state, or the agent's makeup, as input to both, and the reading is a common cause, with nothing running backward. (K2) That blanks read tensed, open as future and stated as past, which the tense-free reading refuses. (K3) That the settling operator's direction, the agent's, is exported to nature. (K4) That the composition primitive is oriented: on the physics face the kernel composes one way and records sit at nodes, stated content of the composition axiom that supplies a mechanism history for physical records, and never the content of (G2). Where all four hold the classical picture is right on its own ground; where (K1) fails the account still runs, and it is at the predictor and the delayed choice that the classical vocabulary produces paradox from a substrate that contains none.

\section{The two horns, and the factorization theorem}
\label{sec:theorem}

\begin{definition}[Behavioral horns]
\label{def:horns}
Fix an attribute stock, the families available to the state; the classification is relative to it. Let $A$ and $B$ be correlated at the mixture level, all attributes of the stock in disagreeing superposition. The arrangement is \emph{mediated} iff there exists a third family $C$ --- elementary or compound, a set of families disjoint from $A$ and $B$, conditioned jointly --- such that conditioning on $C$ screens: within every $C$-pocket, $A$ and $B$ are independent. The arrangement is \emph{direct} iff no such family exists --- that is, for \emph{every} third family, elementary or compound and disjoint from $A$ and $B$, some within-pocket dependence survives its conditioning. The direct horn is the negation of the mediated horn, so the disjunction is exhaustive by construction, and Theorem~\ref{thm:factorization} holds verbatim with $C$ compound. The names that say only what the horns are: $\mathsf{Screenable}(A, B)$ iff some disjoint family $C$ gives $A \perp B \mid C$, and $\mathsf{Unscreenable}(A, B)$ iff for every such $C$, $A \not\perp B \mid C$. ``Mediated'' and ``direct'' are the traditional readings; outside the pairwise sector $\mathsf{Unscreenable}$ is the residual class of dependence no conditioning removes, and it does not mean $A \to B$ or $B \to A$.
\end{definition}

\noindent Screenability is relative to the stock, and refactoring is not adding attributes: for any correlated $P(A, B)$, the attribute $M = (A, B)$ with $P(M = (a, b)) = P(A = a, B = b)$ and hard constraints fixing $A$ and $B$ to its components gives $A \perp B \mid M$, and the representation theorem's auxiliary attribute for a single merged factor is exactly this screener, so an auxiliary screener is not thereby an explanation or a common cause. The horns are defined by conditioning behavior rather than by linkage topology because the two classifications differ: graphical separation implies conditional independence, and the converse fails even for pairwise entries, the three equalities $a = b$, $a = c$, $b = c$ giving an $A$--$B$ factor with $A \perp B \mid C$ all the same. Defining the horns by the test makes the substantive question which jointness realizes which horn; for pairwise-composed material the answer is one displayed line, and for all material it is Theorem~\ref{thm:reichenbach}.

\begin{theorem}[Screening as factorization]
\label{thm:factorization}
Suppose every entry touching the triple $(A, B, C)$ factors pairwise through $C$: the joint weight satisfies $W(a,b,c) = f(a,c)\, g(b,c)$ up to single-family factors. Then for every $C$-pocket $c$ of positive weight, $W(c) = \sum_{a', b'} f(a', c)\, g(b', c) > 0$,
\begin{equation}
\label{eq:factorization}
W(a, b \mid c) \;=\; \frac{f(a,c)\, g(b,c)}{\sum_{a', b'} f(a',c)\, g(b',c)} \;=\; \frac{f(a,c)}{\sum_{a'} f(a',c)} \cdot \frac{g(b,c)}{\sum_{b'} g(b',c)} \;=\; W(a \mid c)\, W(b \mid c),
\end{equation}
the denominator factoring because the double sum is a product of single sums: conditional independence at every positive-weight pocket, and no statement at pockets of weight zero, where conditioning is undefined. Pairwise-composed jointness realizes the mediated horn, and Reichenbach's screening certificate is the factorization signature of pairwise composition.
\end{theorem}

\begin{proof}
Conditioning on $c$ fixes the $C$-argument of every factor and renormalizes; single-family factors in $a$ or $b$ alone are absorbed into $f$ and $g$ respectively; the display is immediate, and independence at fixed $c$ follows from the product form.
\end{proof}

The factorization theorem needs its entries to compose pairwise through one family. The general case needs nothing beyond the multiplicative update of Section~\ref{sec:framework}: a state is a product of entry-factors, and a product carries a separation structure. A \emph{chain} from $A$ to $B$ is a sequence of families beginning at $A$ and ending at $B$, each consecutive pair touched by a common entry; a family $C$, disjoint from $A$ and $B$, \emph{separates} them when every chain from $A$ to $B$ passes through a member of $C$. A state can be written as a product of entries in more than one way, since entries can be merged, and each writing is a \emph{presentation}.

\begin{theorem}[Connectivity, separation, and screening for factorized states]
\label{thm:reichenbach}
Let a state in disagreeing superposition be presented as a product of entries, and let $A$ and $B$ be correlated at the mixture level.
\begin{enumerate}
\item[(i)] \emph{Demand.} A chain of entries joins $A$ and $B$.
\item[(ii)] \emph{Disjunction.} Either some entry touches both $A$ and $B$, or the families other than $A$ that share an entry with $A$ form a family $N_A$ that separates $A$ from $B$.
\item[(iii)] \emph{Certificate.} Every family that separates $A$ from $B$ screens them. A family $C$ screens $A$ from $B$ exactly when the $(A, B, C)$-marginal of the state composes pairwise through $C$, $W(a,b,c) = f(a,c)\, g(b,c)$.
\item[(iv)] \emph{Fork.} Every family $C$ that screens $A$ from $B$ is correlated with $A$ and with $B$, and carries their whole covariance: for all functions $\phi$ of $A$ and $\psi$ of $B$, $\mathrm{Cov}(\phi, \psi) = \mathrm{Cov}\big(\mathbb{E}[\phi \mid C], \mathbb{E}[\psi \mid C]\big)$. When $A$, $B$ and $C$ are binary this is $\mathrm{Cov}(A, B) = P(c)\,P(\bar{c})\,\Delta_A \Delta_B$, with $\Delta_A = P(a \mid c) - P(a \mid \bar{c})$ and $\Delta_B$ likewise, so that for a binary screener of a positively correlated binary pair, labeling $c$ to raise $A$ makes it raise $B$, Reichenbach's fork inequalities.
\end{enumerate}
Parts (i)--(iii) hold in every presentation of the state. In particular, an unscreenable correlation has an entry touching both correlates in every presentation.
\end{theorem}

\noindent Its relation to Reichenbach's principle is this: (i)--(iv) are the statistical consequences associated with the principle, and what they establish are factorization consequences of a factorized state, not causal common ancestry. The fork inequalities are the binary case of (iv) and do not extend to a uniform relevance ordering for compound or multivalued screeners: let $C$ take three equiprobable values with $P(A = 1 \mid C) = 0.1, 0.5, 0.9$ and $P(B = 1 \mid C) = 0.1, 0.9, 0.5$, $A$ and $B$ conditionally independent; then $C$ screens and $\mathrm{Cov}(A, B) = \tfrac{4}{75} > 0$, yet between the last two pockets $A$'s probability rises while $B$'s falls. The generalized common-cause literature distinguishes screening from such additional relevance conditions \cite{hoferszabo2013principle}, and the theorem claims only the covariance identity for every screener and the inequalities for binary ones.

\begin{proof}
(i) If no chain joins $A$ and $B$, every entry touches only families reachable from $A$ or only families that are not, so the weight is a product of a factor in the first families and a factor in the second, $B$ among them; $A$ and $B$ are then independent, against the hypothesis. (ii) If no entry touches both, a chain from $A$ leaves $A$ through an entry that does not touch $B$, and so enters a family of $N_A$. (iii) Let $C$ separate $A$ from $B$, and let $S$ be the families reachable from $A$ by chains avoiding $C$. No entry touches both a family of $S$ and a family outside $S$ and $C$, so the weight is a product of a factor in $S$ and $C$ and a factor in the remaining families and $C$. At each value of $C$ the two factors are functions of disjoint families, which is conditional independence of $S$, and so of $A$, from $B$. For the equivalence, screening gives $W(a,b,c) = W(a,c)\, W(b,c) / W(c)$ at every $c$ of positive weight, which is pairwise composition, and pairwise composition gives screening by Theorem~\ref{thm:factorization}. (iv) If $C$ were independent of $A$, then $A \perp B \mid C$ and $A \perp C$ would give $A \perp B$, against the hypothesis; likewise for $B$. The covariance identity is the law of total covariance with its within-pocket term zero, and for binary $C$ it reduces to the displayed product. The last claim is (ii) and (iii) read in every presentation.
\end{proof}

\noindent Theorem~\ref{thm:reichenbach} gives the statistical screening structure associated with Reichenbach's principle, and two limits belong with it. A separator screens and carries the covariance, but the theorem does not show that the separator is a prior common cause, nor exclude a common effect; that reading needs the mechanism history of Section~\ref{sec:rungs}. And since any joint law can be presented as one factor touching every variable, the demand for a chain has force only relative to a presentation whose factors are the state's stated entries, not an arbitrary refactoring. Part (i) says a correlation is carried by a chain of joint couplings between its correlates; (ii) gives each horn a structural sufficient condition one way, a separating family screening and an unscreenable correlation having a coupling touching both correlates in every presentation, the converse failing, since the $A = B = C$ example is screened by $C$ with an $A$--$B$ factor present; (iii) is the certificate, with the converse the behavioral definition left open; and (iv) returns the covariance identity, with the fork inequalities in the binary case. The steps from product form to screening are the Markov properties of factorized distributions, standard in the graphical-models literature \cite{lauritzen1996graphical}. What differs is where the product form comes from. A causal model obtains it from a Markov condition relative to a graph, with its own arguments for that condition; here it is the multiplicative update of the state, and the statistical consequences follow from that composition without a producer arrow being assumed, the directional information a causal graph carries being located instead in the supplied mechanism history of Section~\ref{sec:rungs}. The theorem was checked in exact arithmetic over 1{,}500 random product-form states on three to five families, with zero weights for joint invalidity and a merged presentation of each: 4{,}208 correlated pairs, 3{,}692 separating families and 2{,}908 screeners, without exception to any part.

Two remarks complete the classification. First, a direct factor $h(a,b)$ that no choice of $f$ and $g$ can absorb survives conditioning on $C$ intact, so that $C$ does not screen; the direct horn needs this for every candidate family, since it quantifies over all of them, and a factor that merely appears in one factorization settles nothing: with $A = B = C$ one fair bit, the weight can be written with $h(a,b) = \mathbf{1}[a = b]$, yet conditioning on $C$ screens. Second, irreducible higher-arity entries fall where the behavioral definition puts them. A parity-style constraint binding $(a,b,c)$ triples jointly, under which conditioning on $C$ forces a direct $A$--$B$ law inside every pocket, touches $A$ and $B$ in one ternary factor, so that $C$ does not separate them; under the uniform weighting $A$ and $B$ are independent before conditioning, so the arrangement belongs to neither horn and is an example of dependence induced by conditioning, the classical collider $A \to C \leftarrow B$, which a causal model represents formally. What the example shows is that topological mediation and behavioral screening come apart.

The triad is then recovered as readings of arrangements. The direct horn admits the two readings, ``$A$ causes $B$'' and ``$B$ causes $A$,'' one arrangement read in either direction. The mediated horn admits three, ``$C$ causes both'' and the two chain tellings, all sharing one screening signature and therefore observationally indistinguishable, their mechanisms and interventions differing as Section~\ref{sec:rungs} shows; the chain and fork tellings are an illustrative class within the screenable horn and not an exhaustive account of it. The three mediated tellings form a Markov equivalence class in the sense of Verma and Pearl \cite{verma1990equivalence, spirtes2000causation, pearl2009causality}, and the present derivation supplies what the equivalence classes presuppose, the screening signature, from coupling composition rather than from graph axioms. The statistical consequences associated with Reichenbach's principle are thereby theorems of the coupling structure in disagreeing superposition, its demand and disjunction by Theorem~\ref{thm:reichenbach}(i)--(ii), its certificate by (iii) and the factorization, and its fork inequalities by (iv) in the binary case, with no production primitive required by the observational results, interventions consuming directed mechanism structure instead.

\section{The Principle Beneath the Theorem}\label{sec:principle}

% >>> ADDED (Fable 5.1, Sep 2 2026, author-directed content enhancement): the positive principle the
% factorization theorem instantiates, stated once and boxed --- the paper had ``no arrow appears''
% as its refrain but no name for what stands where the arrow stood.
The factorization theorem is one instance of a principle that stands where the arrow stood, and the regime boundary of Section~\ref{sec:boundary} is read through it.

\begin{principle}[Joint coupling in place of causation]\label{prin:joint}
A correlation between attributes, relative to an information state, is accounted for not by a mechanism connecting them --- no signal, no shared antecedent, no time order --- but by the \emph{joint coupling}, relative to that state, of the combinations the correlation constrains: the weights the state's entry-factors assign across the combinations, within the admissible structure that joint invalidity fixes, outright joint invalidity being that structural layer and not a limit of the weights. The pruning removes the jointly invalid combinations from the state's standing alternatives, the surviving combinations carry their weights, and the coupling among the survivors \emph{is} the correlation---which it must be, since a correlation need not exclude any combination at all: two binary attributes at weights $0.4, 0.1, 0.1, 0.4$ are strongly correlated with every combination valid. The dependence between $A$ and $B$ is thereby represented by their induced marginal coupling, supported in the specified factorization by a chain of entry-factors, with the joint invalidity of a combination as its deterministic case. A low joint weight on $\{A, \neg B\}$ is not by itself dependence, two independent events of weights $0.001$ and $0.5$ giving $0.0005$ to that cell, and no weight table identifies a direction or an intervention response, which are the mechanism history's (Section~\ref{sec:rungs}); ``$A$ causes $B$'' is reserved for a specified mechanism history and intervention rule.
\end{principle}
\noindent Three consequences follow. First, the classical triad is the inventory of readings of graded coupling in disagreeing superposition: when the coupling can be sorted by a third family whose pockets screen, the marginal factorizes and the common-cause reading is available; when it cannot, the direct reading is, in one of its two directions; the readings differ in the direction read and not in the arrangement read, which is Section~\ref{sec:theorem}'s theorem restated. Second, the principle is regime-neutral where the triad is not: it accounts for a correlation whether or not the alternatives exclude one another, because joint invalidity is a fact about combinations and not about pockets, and that is why it continues to apply in agreeing superposition, where the Bell-local screening representation of Theorem~\ref{thm:reductio} is excluded. Third, no-signaling is characterized at the level of the product by \emph{far-wing balance}: collect every entry-factor involving the far setting or the far outcome into the far wing's joint factor; the restriction is that this product, summed over the far outcome's alternatives, splits into a part free of the far setting and a part free of the near outcome. Proposition~\ref{prop:nosignal} proves that this is exactly what normalized no-signaling needs, that balance of each factor separately is not enough, and the singlet's single coupling factor satisfies it; what explains the balance is a separate question, taken up after the proposition.

\begin{proposition}[No-signaling, at the level of the product]\label{prop:nosignal}
Let the state's weights on the near outcome $a$, far outcome $b$, near setting $\mathbf{x}$ and far setting $\mathbf{y}$ be of the form $W = F(a, \mathbf{x})\, G(a, b, \mathbf{x}, \mathbf{y})$ with $F > 0$, $G \ge 0$ and normalization $Z(\mathbf{x}, \mathbf{y}) = \sum_a F(a, \mathbf{x})\, S(a, \mathbf{x}, \mathbf{y}) > 0$, where $G$ is the product of every entry-factor involving $b$ or $\mathbf{y}$, and write $S(a, \mathbf{x}, \mathbf{y}) = \sum_b G(a, b, \mathbf{x}, \mathbf{y})$. Then the normalized near marginal $P(a \mid \mathbf{x}, \mathbf{y})$ does not depend on $\mathbf{y}$ if and only if $S(a, \mathbf{x}, \mathbf{y}) = g(a, \mathbf{x})\, h(\mathbf{x}, \mathbf{y})$ for some functions $g \ge 0$ and $h > 0$, the converse holding with $g = P(a \mid \mathbf{x}) / F(a, \mathbf{x})$ and $h = Z$. Balance of each factor of $G$ separately does not suffice: with $a \in \{0, 1\}$, $b \in \{-1, 1\}$ and two factors $1 + t\, ab$, where $t = 0$ at one far setting and $t = \tfrac12$ at the other, each factor sums over $b$ to $2$, while their product sums to $2(1 + t^2 a^2)$, and with $F$ uniform the near marginal $P(a = 1)$ is $\tfrac12$ at the first setting and $\tfrac59$ at the second.
\end{proposition}
\begin{proof}
$P(a \mid \mathbf{x}, \mathbf{y}) = F(a, \mathbf{x})\, S(a, \mathbf{x}, \mathbf{y}) \big/ \sum_{a'} F(a', \mathbf{x})\, S(a', \mathbf{x}, \mathbf{y})$. If $S = g h$, the factor $h(\mathbf{x}, \mathbf{y})$ cancels between numerator and denominator. Conversely, if the marginal does not depend on $\mathbf{y}$, then $S(a, \mathbf{x}, \mathbf{y}) = P(a \mid \mathbf{x})\, Z(\mathbf{x}, \mathbf{y}) / F(a, \mathbf{x})$, with $Z$ the denominator, which is the factorization. The counterexample is the displayed arithmetic. Exact checks confirm it, the singlet's balance at two thousand random setting pairs, and the equivalence on three thousand random positive tables.
\end{proof}
\noindent One reading attaches to the principle: relative to any one information state, every combination in that state's block is constrained by the block's existing, so that a correlation is what that constraint looks like from inside, not more determinism than the classical picture but more visibility. The fuller development calls this \emph{predetermination per state} and derives it from the quantifier structure of its pruning rule, which ranges over the whole refinement future of a state at once: with $\mathrm{Base}$ the base contradictions, doom is the least fixed point of $D(A) \iff A \in \mathrm{Base} \,\lor\, \forall B \preceq A\, \exists C \preceq B\, D(C)$ and validity the greatest fixed point of $V(A) \iff A \notin \mathrm{Base} \,\land\, \exists B \preceq A\, \forall C \preceq B\, V(C)$, the recursive equations alone not selecting these solutions \cite{Aethus}. Here $B \preceq A$ is declared proper-child descent, reflexive and transitive, so that both quantifiers range over refinements at every depth; a proper child is a refinement carried at non-negligible relative weight in some admissible decomposition, so that a component of negligible relative weight, zero included, is no proper child and supplies no validity floor. A branch with a safe continuation cone keeps its state valid, and only a future with no such cone dooms the present, with no story required about how the future reaches back.

\noindent\textit{Interpretive restraint.} The balance condition of Proposition~\ref{prop:nosignal} is equivalent to normalized no-signaling, so imposing it characterizes the independence rather than explains it. The explanatory question, which Wood and Spekkens press, is what independently motivated model structure guarantees balance under the allowed changes of settings and mechanisms \cite{wood_spekkens_2015}: a local trace-preserving quantum construction supplies such a guarantee, an unrestricted factor model constrained to balance does not, and an arrowless nonnegative coupling model can signal, as the perturbation in Section~\ref{sec:objections} shows. The proposition's nonnegative hypotheses cover the singlet directly, $F = 1$, $G = 1 - ab\,\mathbf{x}\cdot\mathbf{y}$, $S = 2$, $Z = 4$, its zeros at aligned and anti-aligned settings included. In this paper, therefore, the statement that no coupling value is tuned is attached to the specified quantum construction and its permitted local operations, imported from the quantum companion \cite{AethusQuant}, and not to the arrowless form of the coupling or to the imposition of balance.

\section{Two instances, one of them double}
\label{sec:instances}

One instance in the literature runs the certificate twice in one vignette. In the smoking-lesion case of the medical Newcomb family \cite{skyrms1980causal, nozick1969newcomb}, a lesion is stipulated as the common cause of cancer and of a taste for smoking, smoking itself causally inert. Conditioning on the lesion screens smoking from cancer --- the vignette's inertness stipulation \emph{is} Theorem~\ref{thm:factorization} with $C$ the lesion. And, under the routing assumption the decision companion states, that the taste is the only channel from the lesion to the act so that the act is independent of the lesion given the taste, which causal inertness alone does not supply, conditioning on the stated taste screens the agent's act from the lesion --- which is the tickle defense of the evidential-decision-theory literature \cite{eells1982rational}, here appearing as Reichenbach's screening certificate applied at the mediator: a decided pocket leaves the mixture-level averaging nothing further to move. The certificate structure itself is no news to the tickle literature, whose apparatus is explicitly Reichenbachian \cite{eells1982rational}; what is claimed is narrower, the certificate obtained from coupling composition and the two-certificates-in-one-vignette count; the decision-facing consequences --- which stated couplings bind on an act's settling, and how the criterion sorts Newcomb, the lesion, and their transparent variants with one rule --- are developed in the decision extraction \cite{AethusDecision} and not re-argued here.

\section{The interventional and counterfactual rungs}
\label{sec:rungs}

Pearl's hierarchy has three rungs, seeing, doing and imagining \cite{pearl2009causality}. The paper's contribution to them runs in four steps, graded accordingly: separation and screening in product presentations (Section~\ref{sec:theorem}); the equivalence of settling along a mechanism history with intervention in a causal Bayesian network (Theorem~\ref{thm:settling}); the additional structure, functional entries over background attributes, that counterfactual equivalence needs (Proposition~\ref{prop:cf}); and the philosophical reading of those structures (Section~\ref{sec:concept}). Theorem~\ref{thm:reichenbach} is the first, and conditioning itself is the framework's pruning followed by renormalization \cite{Aethus}. This section takes the other two, and it begins with the one fact about them that holds for every theory.

\begin{proposition}[No theory reads interventions off couplings alone]
\label{prop:parity}
There are states with one joint coupling and different interventional laws. With $U$ a fair bit, the state in which $A = U$ and $B$ copies $A$, and the state in which $B = U$ and $A$ copies $B$, carry the same joint weight, $W(0,0) = W(1,1) = \tfrac12$; setting $A$ to $0$ fixes $B = 0$ in the first and leaves $B$ uniform in the second.
\end{proposition}
\begin{proof}
Both states put weight $\tfrac12$ on each agreeing pair and none elsewhere. Setting $A$ replaces $A$'s own entry in each: in the first, $B$'s entry still copies $A$; in the second, $B$'s entry is its fair coin.
\end{proof}

\noindent A causal Bayesian network takes its graph as input for exactly this reason, and the proposition applies to it equally: no account, with arrows or without, recovers interventional laws from the joint coupling alone. What an account can do is say where the further structure lives. A causal model puts it in a posited graph; the framework puts it in a state's mechanism history, a structure supplied beside the state and distinct from the order of its observational updates, which is acquisition history and carries no direction. Call a \emph{mechanism history} a sequence of entries $e_1, \ldots, e_n$ in which $e_k$ touches a family $V_k$ and some earlier families, its \emph{inputs} $\mathrm{in}_k$, and sums to one over $V_k$ at every value of its inputs; $e_k$ \emph{establishes} $V_k$. In the \emph{default binding}, settling a family suspends the entry that established it and nothing else: that entry is how the family's value was produced, and every other entry states how the rest depends on it. This is the decision extraction's settling operator $\mathrm{Settle}_\mu$, with its binding modes read off the history \cite{AethusDecision}.

\begin{theorem}[Settling is surgery]\label{thm:surgery}
\label{thm:settling}
Let a state be given by a mechanism history. Its weight is $W(v) = \prod_k e_k(v_k \mid \mathrm{in}_k)$, and settling $V_k$ to $x$ in the default binding gives
\begin{equation}
\label{eq:truncated}
W_{V_k = x}(v) \;=\; \mathbf{1}[v_k = x] \prod_{j \ne k} e_j(v_j \mid \mathrm{in}_j),
\end{equation}
which is Pearl's truncated factorization for the causal Bayesian network whose parent sets are the entries' inputs. Conversely, every causal Bayesian network is a mechanism history: list its variables in a topological order and take its conditional tables as the entries.
\end{theorem}
\begin{proof}
The history's product is the network's factorization. Settling multiplies the indicator of $V_k = x$ into the state and suspends $e_k$, which leaves \eqref{eq:truncated}. It is already normalized: summing out the families from the last to the first removes the remaining entries one at a time, each summing to one over its own family, and the indicator sums to one over $V_k$. The converse is the definition of a causal Bayesian network read as a history.
\end{proof}

\begin{corollary}[The calculus is inherited]
\label{cor:inherited}
Every result that causal Bayesian networks derive from the truncated factorization and the Markov property holds for settling on mechanism histories without further argument: the three rules of the do-calculus, and with them the back-door and front-door identification criteria \cite{pearl2009causality}. And the three graphs that share one screening structure, $A \to C \to B$, $A \leftarrow C \to B$ and $A \leftarrow C \leftarrow B$, are three histories of one product whose settlings differ: settling $A$ can move $B$ only in the first, settling $B$ can move $A$ only in the third, and in the second neither can move the other.
\end{corollary}
\begin{proof}
Theorem~\ref{thm:settling} identifies settling on a mechanism history with the network's intervention, so each network theorem applies with its terms read through that identification. The three histories list $A, C, B$, then $C, A, B$, then $B, C, A$, their entries being the conditional tables of one joint law.
\end{proof}

\noindent The orientation a causal graph posits is here the mechanism history's, which mechanisms take which families as inputs, and it is consumed rather than derived: the order in which the state acquired its families, which the refinement history records in any case, is acquisition order and carries no direction (Section~\ref{sec:rerendition}). The fuller development's counterfactual operator reads the mechanism history where one is supplied. Its Mode~P perturbation ascends to the nearest state in which the target is blank, restates it, and re-prunes; that operation is intended to implement Theorem~\ref{thm:settling}'s surgery, but the identification is not automatic, since moving to an earlier state can lose unrelated information while surgery replaces one mechanism and preserves the others, and an equivalence under stated preservation conditions is not proved here; the ascent is indexed by blankness, not by tense \cite{Aethus}. The third rung needs more than the tables, as the next proposition shows, and follows once the mechanism entries are functional.

\begin{proposition}[Counterfactuals need functional structure]\label{prop:cf}
(i) Conditional tables do not fix counterfactuals. With $X$ and $U$ independent fair bits, the mechanism histories $Y = U$ and $Y = X \oplus U$ have the same observational law on $(X, Y)$, though not on $(X, U, Y)$, and the same interventional law $P(Y \mid do(X = x))$ for each $x$, yet given $X = 0$ and $Y = 0$ the first gives $Y_{X \leftarrow 1} = 0$ and the second $Y_{X \leftarrow 1} = 1$.
(ii) With functional structure the correspondence holds. Suppose every mechanism entry is deterministic, $v_k = f_k(\mathrm{in}_k, u_k)$, with the background attributes $u_k$ carried as the state's own attributes under prior entries, and let the evidence $e$ have positive probability. Define the counterfactual settling of $X = x'$ given $e$: use $e$ to infer the posterior $P(u \mid e)$ on the background attributes, the only use made of the factual evidence; settle $X = x'$ by replacing $X$'s entry, as in Theorem~\ref{thm:settling}, retaining that posterior; and recompute every descendant of $X$ through its entry from the background values, the factual values of those descendants not being retained as constraints on the recomputed ones. The resulting weight on $Y = y$ is $\sum_u P(u \mid e)\,\mathbf{1}[Y_{x'}(u) = y]$, the structural causal model's abduction, action and prediction.
\end{proposition}
\begin{proof}
(i) In both histories $P(Y = y \mid X = x) = \tfrac12$, and $X$ having no inputs, $P(Y \mid do(X = x)) = P(Y \mid X = x)$; the evidence fixes $U = 0$ in both, whence the two counterfactual values. (ii) The state is the product of the background attributes' prior entries and the deterministic entries $\mathbf{1}[v_k = f_k(\mathrm{in}_k, u_k)]$. Pruning by the evidence leaves the weights $P(u \mid e)$ on the consistent background combinations, which is abduction, the normalization existing because $P(e) > 0$; replacing $X$'s entry by $X = x'$ is the action; and recomputing each descendant through its deterministic entry is prediction, each background combination carrying exactly one value of $Y$. A factual value of a descendant is a consequence of the background values under the factual mechanism, and it is recomputed, not retained: retaining it would constrain the counterfactual to agree with the actual on a value the surgery changes. An exact check against a brute-force twin-world computation on three thousand random models agrees in every case.
\end{proof}
\noindent So settling establishes intervention equivalence, and the counterfactual rung is recovered when the mechanism history carries functional entries over background attributes, the structure a structural causal model adds to a causal Bayesian network for the same purpose. Whether any of the companion semantics' scope settings performs the counterfactual settling is a question for that semantics \cite{Aethus}; the correspondence is proved here for the operation as defined. Settling validates Cautious Monotony: settling a held consequence at the value it already takes changes no input anywhere, which for structural models is the composition property \cite{galles1998axiomatic}; the companion semantics' \emph{reopening} does not \cite{AethusSemantics}, since it keeps every law and reopens the grounds of a supposed entry, so a held consequence restated can lose a conclusion that leaned on them. The framework carries both conditionals, and Cautious Monotony separates them.

\noindent Three extensions mark where the framework goes past the classical theory rather than recovering it. The default binding is one binding assignment among others: the decision extraction's Newcomb predictor is an entry acquired before the act yet stated robust to the act's mode of production. In the conventional mechanism history the prediction is an input of the act, $\mathrm{Prediction} \to \mathrm{Act}$, and surgery on the act cuts that incoming relation, so the prediction no longer varies with the act and causal decision theory two-boxes; settling that honors the stated robustness retains the binding, the prediction tracking the act, and one-boxes \cite{AethusDecision}. A causal Bayesian network can express the retained relation only with the arrow reversed, $\mathrm{Act} \to \mathrm{Prediction}$, under which surgery on the act keeps the downstream prediction mechanism and $P(\mathrm{Prediction} \mid do(\mathrm{Act} = a)) = P(\mathrm{Prediction} \mid \mathrm{Act} = a)$, as Theorem~\ref{thm:settling} gives; what the framework declines is the reading of that retained binding as causal production of the prediction by the act, the binding being a constraint between entries and not a mechanism. The Bell correlations are the second extension, grounded by graded coupling, with no-signaling characterized by far-wing balance and explained by the imported quantum construction under its specified operations (Section~\ref{sec:boundary}). The third is the counterfactual semantics' Cautious Monotony verdicts, which the fuller development proves no smooth preferential representation reproduces on all of the operator's evaluable verdicts, under its dual-peg reading of conjunctive antecedents \cite{Aethus}. Classical causal theory is thereby the framework's default-binding case in the exclusive regime: recovered at the first two rungs, and at the third given functional mechanism entries. The identification was checked in exact arithmetic over 1{,}500 random four-family networks: settling equals the truncated factorization in every case, and the back-door adjustment over the target's inputs holds wherever it is defined.

\section{The regime boundary} \label{sec:boundary} What this section locates is an obstruction to a specified global valuation extension, and not a failure of Theorem~\ref{thm:factorization}: for performed settings and recorded outcomes the singlet's statistics are a finite classical table, $P(a, b, \mathbf{x}, \mathbf{y}) = P(\mathbf{x}, \mathbf{y})\,(1 - ab\,\mathbf{x}\cdot\mathbf{y})/4$, presentable as a product of factors, to which connectivity and separation apply; what fails is the extension to simultaneous setting-independent values $A_0, A_1, B_0, B_1$, and ``boundary'' below names that extension's failure. Theorem~\ref{thm:factorization} quantifies over pockets, and that is the sense in which the boundary is its. In agreeing superposition the alternatives stand together rather than excluding one another, and the interference phenomenology the co-obtaining licenses is what the exclusive arrangement cannot host. A correlation grounded at the co-obtaining grade, the Bell correlations being the standing example, the entangled state analyzed in the fuller development as a co-obtaining superposition across the two sites \cite{Aethus}, satisfies the grounding demand in full, being carried by joint couplings as ever, while escaping the mediated horn's certificate in one specified sense: no mediator inducing a representation that satisfies (D1)--(D3) of Theorem~\ref{thm:reductio} exists, which in the standard scenario is a setting-independent $C$ with local responses, $P(a, b \mid \mathbf{x}, \mathbf{y}, C) = P(a \mid \mathbf{x}, C)\,P(b \mid \mathbf{y}, C)$ and $P(C \mid \mathbf{x}, \mathbf{y}) = P(C)$. Screening alone is weaker than this and is not excluded, as the instance below shows, and the auxiliary $M = (A, B)$ of the screenability remark screens the observed outcomes trivially without supplying a grounding of the required kind. Bell violation obstructs the specified global, measurement-independent valuation model; it does not invalidate ordinary conditional independence or the factorization results for the observed probability table.

The relation to the fine-tuning result is this. Wood and Spekkens showed that a classical causal \emph{model} of the Bell statistics must tune its parameters so that posited causal pathways leave no marginal signature \cite{wood_spekkens_2015}. The present structure is not a classical causal model and posits no pathways; that does not by itself discharge the fine-tuning burden, since an arrowless nonnegative coupling model can signal. The marginals are characterized instead by far-wing balance, a structural restriction on the far wing's joint factor (Proposition~\ref{prop:nosignal}), and what guarantees the balance under the permitted operations is the imported quantum construction, whose two-tag form obtains no-signaling as an algebraic identity \cite{AethusQuant}. The levels are distinct: the \emph{outcome statistics} are pocketed, stated records over which conditional independences can be asked and answered classically, while the obstruction lives at the \emph{grounding} level, where the correlation's carrier is co-obtaining and no (D1)--(D3) mediator exists.

What a \emph{setting choice} is in this framework, since the fine-tuning theorem concerns causal models with settings as free variables: a setting is an act attribute, blank until settled and stated thereafter; its outcomes and records live at the pocketed level, where the balance identity governs their statistics; and what makes a setting free, which stated couplings bind on its settling and under which modes of production, is the binding-modes analysis of the decision extraction, developed there and not here \cite{AethusDecision}. The obstruction concerns the correlation's grounding, not the settings' freedom.

The section's thesis, at the strength its theorem supports: the Bell-type deviation from Reichenbachian explanation is the boundary of a specified representation, the (D1)--(D3) valuation extension, and not an exotic causal species. The same primitive that yields the screening structure in disagreeing superposition grounds the correlations past that boundary; the screening structure associated with Reichenbach's principle is the coupling's exclusive-regime specialization, and the Bell correlations are its co-obtaining face.

\subsection{The boundary as a theorem}\label{sec:boundary_theorem}

% >>> ADDED (Fable 5.1, Sep 2 2026, content enhancement): the regime boundary, which the section above
% establishes by the quantification losing its subject, now also established as a theorem --- the
% companion's Deflation Reductio in neutral vocabulary --- with the vacuum reading of why the co-obtaining
% correlations are unmediated and a Bell instance worked in the paper's terms.
The section above locates the boundary at the failure of a (D1)--(D3) mediator. The theorem states that failure exactly: no reparsing of the regime into exclusive pockets of the (D1)--(D3) kind reproduces Bell-violating statistics.

\begin{theorem}[Deflation reductio, boundary form]\label{thm:reductio}
Call an assignment $\delta$ a \emph{lossless deflation} of an information state $A$ when it provides, for $A$ and every state in $A$'s refinement structure, a family $W$ of complete valuations of the attributes such that: (D1) each valuation assigns every attribute one value, independent of which other attributes are stated or measured, agreeing with the state's stated entries; (D2) on the exclusive fragment, a state is valid iff its family is nonempty, with refinement restricting the family; (D3) the ensemble statistics over any attributes of $A$ are the marginals of one probability measure on $W(A)$, each run's outcomes identified with that run's values, the choice of which attributes are measured being independent of the valuation. Then: (a) if the blank attributes of $A$ contain a Bell quadruple --- four attributes whose CHSH combination exceeds $2$ --- no lossless deflation of $A$ exists; (b) in the Bell scenario proper---two parties, two settings each, two outcomes---if the full family of CHSH inequalities holds together with the compatibility of the marginals, Fine's theorem supplies a joint probability distribution over the four attributes; whether that distribution can be organized into a family $W$ over the whole refinement structure satisfying (D1)--(D3) is a further step, and (b) claims only the existence of the joint distribution on that scenario, not a lossless deflation of an arbitrary state.
\end{theorem}
\noindent Part (a) is the theorem; (b) is its converse on the one scenario where the converse is known, stated at the strength Fine's result supports. A stronger (b), claiming a lossless classical filtration for any collection of blank attributes satisfying CHSH, would not follow: CHSH is complete for the $2 \times 2 \times 2$ scenario only, and even there the joint distribution Fine supplies must still be shown to extend across refinements with (D1)--(D3) intact.
\noindent\textit{Proof.} (a) Suppose $\delta$ exists. By (D1) each valuation assigns a value to all four attributes, so the measure of (D3) lives on the four-fold value space, and by (D3) the four pairwise correlations are its marginals, measurement independence letting one measure serve all four setting pairs. By Fine's theorem the existence of such a joint distribution is equivalent to the CHSH inequalities \cite{fine1982hidden, pitowsky1994george}, so the combination is at most $2$, contradicting the hypothesis. (b) On the Bell scenario, Fine's theorem in its converse direction supplies a joint probability distribution over the four attributes from the full CHSH family together with marginal compatibility \cite{fine1982hidden}; that is the whole of what (b) asserts. The further construction a stronger (b) would need---taking $W$ to be the complete valuations consistent with each state's closed content with the measure pushed forward, so that (D1) holds by construction and (D2) because validity on the exclusive fragment coincides with the existence of a consistent completion---would additionally require that the Fine distribution extend coherently across every refinement of $A$ with (D3) intact, which is not established here and is booked as open. \hfill$\square$

\noindent To read a co-obtaining blank as a set of exclusive pockets whose values are merely unknown is to assert (D1) and (D3); (a) refutes that assertion wherever the correlations violate the bounds, and (b) says that where they do not, on the Bell scenario, the particular obstruction (a) identifies is absent and a joint distribution exists, which is less than a full exclusive-pocket representation across the refinement structure, left open. The obstruction falls at the Bell bounds and is not drawn by hand, but it is an obstruction's edge, not a proof of representability on the far side. Two predicates are distinct: $\mathsf{CoObt}(A)$, the framework's structural regime, that the alternatives co-obtain at $A$; and $\mathsf{BellViol}(P)$, a statistical property of an observed distribution. What is proved is $\mathsf{BellViol}(P_A) \Rightarrow \neg\,\mathsf{Deflatable}_{\mathrm{D1\text{--}D3}}(A)$; the biconditional $\mathsf{CoObt}(A) \iff \mathsf{BellViol}(P_A)$ is not asserted, a co-obtaining state possibly generating Bell-local statistics at some settings.

\paragraph{The record restriction, stated as a premise.} The theorem excludes a (D1)--(D3) representation and says nothing about why the two sites' alternatives carry no such mediator. The framework's answer is a premise about records, imported from the fuller development and stated here as what it is \cite{Aethus}: a co-obtaining superposition admits no record, in any information state of the block, past or future, that reveals a prior selection among its alternatives, a context-independent value already there. Were such a record possible, the alternatives would already exclude one another relative to the state carrying it, and the superposition would be exclusive rather than co-obtaining; an eventual measurement records a newly settled outcome in its context, and that record is compatible with the superposition. The restriction is on information access, not a mechanism of concealment: nothing holds the distinction in, and the cut between the regimes falls along the line ``no state in the block accesses the alternatives' distinction'' on the co-obtaining side and ``otherwise'' on the exclusive side. The bridge from the theorem's latent variable to recorded content is a second premise, stated here: a determinate mediator admissible in the framework is carried as indexed content of some state of the block, a mathematical latent variable being no physical record by itself. Under it, a mediator of the (D1)--(D3) kind is content that reveals a prior, context-independent selection among the alternatives, which is exactly what the record discipline excludes on the co-obtaining side; so none exists there under the two premises. What is excluded is that packaging of a mediator; contextual, measurement-dependent or nonclassical mediators, the quantum common causes of the causal-model literature among them \cite{allen2017quantum}, are not addressed by the argument, and the claim is not that no mediator of any conceivable kind exists. The statistical exclusion of the (D1)--(D3) class, stochastic models included, is Bell's theorem's, by Fine's equivalence, and needs no premise of the framework's; what the record discipline adds is the framework's explanation of why no admissible state of the block carries such a mediator, a restriction on admissible states and not a further exclusion of models.

\paragraph{A Bell instance, in the paper's terms.} Take the singlet, with Alice's and Bob's outcomes at their respective settings the four attributes of the quadruple. Relative to either experimenter's state before measurement, all four are blank, and the pair of outcomes at any one setting pair is a co-obtaining superposition across the two sites: the combinations \emph{same outcome} and \emph{opposite outcome} are both standing alternatives when $-1 < \mathbf{x}\cdot\mathbf{y} < 1$, with no record anywhere that separates them, and their relative weights are what the entry-factors assign, the anticorrelation at equal settings and the $\cos$-law elsewhere. What each measurement does is re-index rather than act: the near outcome becomes stated content of the near experimenter's state, and the far outcome's weights are conditioned relative to that state by the coupling, $P(b \mid a, \mathbf{x}, \mathbf{y}) = (1 - ab\,\mathbf{x}\cdot\mathbf{y})/2$ for outcomes $a, b \in \{-1, 1\}$, so that the far outcome is predicted with probability one at aligned or anti-aligned settings and is unbiased at perpendicular ones, with nothing traveling. Two distinctions the paper has drawn apply here. The factor $1 - ab\,\mathbf{x}\cdot\mathbf{y}$ assigns zero weight to the same-outcome combination at aligned settings and to the opposite-outcome combination at anti-aligned ones, and at generic settings gives every combination positive weight while remaining correlated, as the positive-table example of Section~\ref{sec:rerendition} says it must; whether the zero-weight recorded combinations are also structurally invalid is not settled by the factor, which lives in the weight layer, but by the quantum companion's record discipline, and the paper does not assert it here. Likewise, probability-one prediction of the far outcome is not its determinacy relative to the near state; determinacy arrives with a record, by the validity rule, and the two are distinguished as they are throughout. As for mediation: a setting-independent screener of the observed outcomes in every context is \emph{not} excluded by the statistics. Let $\lambda \in \{-1, 1\}$ be fair and independent of the settings, let $a = \lambda$, and let $P(b \mid \mathbf{x}, \mathbf{y}, \lambda) = (1 + b\lambda E_{\mathbf{x}\mathbf{y}})/2$ with $E_{\mathbf{x}\mathbf{y}} = (-1)^{xy}/\sqrt2$ for binary settings; then $a$ and $b$ factorize given $\lambda$ in every context, the marginals are unbiased and non-signaling, and the CHSH combination is $2\sqrt2$. What that screener lacks is locality of the response, Bob's conditional law depending on Alice's setting, so it is not a (D1)--(D3) mediator, and it is such mediators, and only such, that Theorem~\ref{thm:reductio} excludes. The correlation is carried by the graded coupling across the two sites, no horn of the triad reads it, and its no-signaling is the far-wing balance the imported construction protects under its specified operations.

\section{The concept, component by component}
\label{sec:concept}

\noindent What follows is a proposal about how the components of the concept could be re-read in the framework's terms. It is not a demonstration that the concept has been replaced, which the mathematics of this paper does not establish.


The results so far are organized by rung, and the concept of causality is larger than its rungs. It bundles an account of production, an account of regularity, a probabilistic account, a principle of common cause, a picture of physical process, an arrow, and an interventionist and a counterfactual semantics, together with the practices built on them: causal discovery, decision, and explanation, and since Bell a verdict on quantum nonlocality. A replacement for the concept owes an account of each. Table~\ref{tab:concept} gives it. Each entry is a theorem of this paper, a reading of Principle~\ref{prin:joint}, or a result of a companion, as marked, and the entries new here are proved or derived below the table.

\begin{table}[!htbp]
\centering\small
\renewcommand{\arraystretch}{1.25}
\begin{tabular}{p{0.17\textwidth}p{0.31\textwidth}p{0.42\textwidth}}
\hline
\textbf{Component} & \textbf{Classical carrier} & \textbf{Replacement here} \\
\hline
Production & a producer relation, the cause bringing the effect about & joint coupling, with joint invalidity its structural layer (Principle~\ref{prin:joint}) \\
Regularity & constant conjunction \cite{hume1748enquiry}; INUS conditions \cite{mackie1965causes} & the minimal jointly invalid sets, each read $k$ ways (Proposition~\ref{prop:regularity}) \\
Probability raising & $P(E \mid C) > P(E)$ \cite{suppes1970probabilistic, eells1991probabilistic} & the mixture-level face of within-pocket coupling, Simpson reversals being between-pocket terms (below) \\
Common cause & Reichenbach's demand, disjunction and certificate & Theorem~\ref{thm:reichenbach} \\
Physical process & transmission of a mark or a conserved quantity \cite{salmon1984scientific, dowe2000physical} & a chain of entries, with nothing travelling along it (Theorem~\ref{thm:reichenbach}(i); below) \\
Direction & causes precede their effects & three orders separated: mechanism input order (which entries establish which families), binding, and time of acquisition (below) \\
Intervention & graph surgery \cite{pearl2009causality}; invariance under ideal intervention \cite{woodward2003making} & settling along a supplied mechanism history, invariance as binding (Theorem~\ref{thm:settling}) \\
Counterfactual dependence & similarity orderings \cite{lewis1973causation}; structural equations \cite{halpern2016actual} & counterfactual settling on functional mechanism entries, Pearl's abduction--action--prediction (Proposition~\ref{prop:cf}); its correspondence with the companion's reopening operator left open \cite{AethusSemantics} \\
Discovery & Markov equivalence classes \cite{verma1990equivalence, spirtes2000causation} & the screening signature derived, and the class as histories of one product (Corollary~\ref{cor:inherited}); what observation cannot fix, no theory can (Proposition~\ref{prop:parity}) \\
Decision & act-to-outcome arrows & binding on the settling: Newcomb one-boxes, the lesion smokes, transparent Newcomb two-boxes with positive predictor error, the companion's admissibility rule governing the perfect endpoint \cite{AethusDecision} \\
Explanation & citing causes & citing the coupling, the triad being its inventory of readings (Principle~\ref{prin:joint}) \\
Quantum nonlocality & fine-tuned classical causal models \cite{wood_spekkens_2015} & the regime boundary: grounded by coupling, the (D1)--(D3) valuation extension excluded (Theorem~\ref{thm:reductio}); no-signaling as far-wing balance (Proposition~\ref{prop:nosignal}), explained by the imported quantum construction \cite{AethusQuant} \\
\hline
\end{tabular}
\caption{The classical concept of causality, component by component, and the proposed counterpart of each component here.}
\label{tab:concept}
\end{table}

\begin{proposition}[Regularity is the structural layer]
\label{prop:regularity}
Let a valid state have only deterministic entries, so that its content is fixed by which combinations are jointly invalid, and let its attributes be binary, $\bar{x}$ denoting the other value of $x$. A set of values is \emph{jointly invalid} when no valid combination contains it, and \emph{consistent} otherwise.
\begin{enumerate}
\item[(i)] A consistent set $S$, possibly empty, is a minimal sufficient condition for $e$---every valid combination containing $S$ contains $e$, and no proper subset of $S$ has this property---exactly when $S \cup \{\bar{e}\}$ is a minimal jointly invalid set.
\item[(ii)] So $c$ is an INUS condition for $e$ in Mackie's sense exactly when $c$ belongs to a set $S$ of at least two values such that $S \cup \{\bar{e}\}$ is minimal jointly invalid and some valid combination contains $e$ without containing $S$.
\item[(iii)] Each minimal jointly invalid set $\{x_1, \ldots, x_k\}$ with $k \ge 2$ is read $k$ ways: for each $j$, the other members form a minimal sufficient condition for $\bar{x}_j$.
\end{enumerate}
\end{proposition}
\begin{proof}
(i) Every valid combination containing $S$ contains $e$ exactly when none contains $S \cup \{\bar{e}\}$, that is, when $S \cup \{\bar{e}\}$ is jointly invalid. A proper subset $T$ of $S$ is sufficient for $e$ exactly when $T \cup \{\bar{e}\}$ is jointly invalid, and the proper subsets of $S \cup \{\bar{e}\}$ are these sets together with $S$ itself, which is consistent; so minimality corresponds on the two sides. (ii) Mackie asks that $c$ be a non-redundant part of a condition that is sufficient but not necessary, $c$ alone not sufficing. By (i) the sufficient condition is $S$; its non-redundancy is minimality; its not being necessary is a valid combination containing $e$ without $S$; and $c$'s insufficiency alone is $|S| \ge 2$. (iii) Each $\{x_i : i \ne j\}$ is a proper subset of a minimal jointly invalid set, hence consistent, and adding $x_j$ to it gives that set; by (i) with $e = \bar{x}_j$, it is a minimal sufficient condition for $\bar{x}_j$.
\end{proof}

\noindent Regularity theories, from Hume's constant conjunction to Mackie's INUS conditions \cite{hume1748enquiry, mackie1965causes}, thereby correspond to the structural layer of joint coupling, read through its minimal jointly invalid sets, and both of their classic failures are visible in the proposition. Part~(iii) is the direction problem. A minimal jointly invalid set carries no mark of which member is the effect, so regularity alone cannot tell the flagpole's height from its shadow's length as the explaining term; the orientation comes from the mechanism input structure the settling theorem assumes (Theorem~\ref{thm:settling}), a topological ordering recording part of that structure without replacing it, not from the set. And a regular conjunction of two effects of one cause---the barometer's fall and the storm, or Mackie's own Manchester hooters and the London workers leaving at five \cite{mackie1974cement}---is screened by the family that separates them (Theorem~\ref{thm:reichenbach}(iii)). The proposition was checked by brute force over 600 random deterministic states on four binary attributes, $118{,}450$ condition--outcome pairs in all, without exception.

\noindent Probability raising, the core of the probabilistic theories \cite{suppes1970probabilistic, eells1991probabilistic}, is a mixture-level quantity, and the law of total covariance splits it along any family $C$:
\begin{equation}
\label{eq:total-covariance}
\mathrm{Cov}(A, B) \;=\; \mathbb{E}\big[\mathrm{Cov}(A, B \mid C)\big] \;+\; \mathrm{Cov}\big(\mathbb{E}[A \mid C],\, \mathbb{E}[B \mid C]\big).
\end{equation}
The first term is the coupling within $C$'s pockets, the second the coupling carried between them. A Simpson reversal \cite{simpson1951interpretation}---raising in the mixture, lowering in every pocket---is the second term outweighing and opposing the first, and Cartwright's requirement that a causal law raise the probability of its effect in every relevant context \cite{cartwright1979causal} is the requirement to read the coupling within pockets. Theorem~\ref{thm:reichenbach}(iv) is the case in which the within-pocket term vanishes and the whole correlation is carried between pockets.

\noindent Process theories \cite{salmon1984scientific, dowe2000physical} locate causation in a continuous process transmitting a mark or a conserved quantity. The chain of entries in Theorem~\ref{thm:reichenbach}(i) is such a process's skeleton with the transmission removed: a sequence of couplings, consecutive ones sharing a family. What the replacement drops is the claim that something travels along the chain to produce the correlation. The chain is the correlation's ground, and where it is realized along a worldline, the process theorist's continuity is that realization.

\noindent The arrow is where the classical concept bundled most: precedence records the order in which things were settled, fixes which dependences survive an intervention, and aligns both with the order of time. The replacement separates the three. The first is the mechanism history supplied with the state, which settling reads (Theorem~\ref{thm:settling}); it is an order of mechanisms fixed by the input sets, which determine a directed acyclic graph, and not an order of inquiry, since learning of smoke before fire leaves the fire among the smoke's inputs. The second is the binding modes, which default to that order and can depart from it, as the Newcomb predictor's robustness does. The third is proper time, which the validity layer does not consult: delayed-choice arrangements, in which the supplied input structure runs against the temporal order of acquisition, are handled without retrocausal edges \cite{Aethus}. The direction is relocated into the input sets, not eliminated, and the time order was never the one doing the causal work.

\noindent Read down the table, the classical concept is the special case the replacements reduce to where binding is default, alternatives exclude one another, and weights are classical. In that sector production reads as coupling, regularity as coupling's structural layer, and the arrow as the mechanism history's direction, and the specified Bayesian-network and functional-model results transfer under the correspondence. Outside it the replacements continue where the concept stops: at Newcomb, at Bell, and at the counterfactual verdicts no ordering follows. Whether that asymmetry makes the proposal a replacement of the concept or a reinterpretation of it is left open here; the ontological claim the paper does make, that an indexed state describes reality relative to its index completely and that its hard structure is not read off its table (Example~\ref{ex:twostates}), is what the replacement reading would rest on.

\section{Placement of the proposal}
\label{sec:dissolution}

The observational analysis consulted no arrow, and the triad's readings did no work in it, being readings of arrangements; the arrows entered only with the mechanism histories of Section~\ref{sec:rungs}, as consumed structure. On this reading the classical insistence on producer categories is a schema built for the pairwise-composed special case of disagreeing superposition and exported beyond it, and the strangeness of the quantum correlations is the export failing where the (D1)--(D3) representation does. The common-cause-system program, partitions of generalized screeners in place of a single screening event \cite{hoferszabo2013principle}, generalizes the certificate while keeping the exclusive regime's furniture, and the present analysis places both its successes and its Bell-front limit: within the regime the certificate is a theorem with room to generalize; at the boundary no enlargement of the screener family satisfying (D1)--(D3) reproduces the Bell-violating statistics, which is the exact form of the limit, while an exclusive description within each measurement context remains available and asserts no single noncontextual, measurement-independent valuation across contexts. A realist account of the kind Spekkens has said would suffice \cite{spekkens2024escapedsapiens} is then available with the mystery's source withdrawn rather than reinterpreted: the epistemic--ontic scramble Jaynes named \cite{jaynes1990probability}, which Schmid, Selby and Spekkens have pressed the field to undo \cite{schmid2020unscrambling}, is unscrambled by construction here, since what is index-relative and what is determinate are never mixed. That the result is a dissolution of the quantum-causation problem rather than a representation of it is the philosophical proposal of Section~\ref{sec:concept}, argued there and not established by the theorems.

\section{Objections}
\label{sec:objections}

\emph{This is the retrocausal program under another name.} The kinship is real: the retrocausal and all-at-once programs \cite{wharton2015universe, price2012does, sutherland2008causally}, and, doing something different, since her constraints govern the whole history without a preferred direction rather than boundary conditions with backward influence, Adlam's global-constraint account of law \cite{adlam2022laws}, all run constraints on a block rather than dynamics through one, and so does the present account, whose substrate is joint invalidities symmetric in time and on which nothing happens on a slice. The difference is the first postulate of the companion works. Those programs constrain one determinate history, or a fixed pair of boundary conditions; here there is no un-indexed content for a constraint to fix, the constraints being validity constraints on information states, and what they carve is which indexed realities can stand together, not which single history obtains. Reduced-form coupling is therefore not a retrocause any more than it is a common cause: both are dynamics-shaped stories about one history, and the pruning is neither.

\emph{Behavioral horns are causal dependence renamed.} The Newcomb prediction entry of the decision extraction is the existence proof to the contrary: an entry can be robust to how the act is produced, binding on the settling, while no act-to-prediction causation exists or is wanted, so binding behavior and causal production come apart in the founding case \cite{AethusDecision}. Behavior under conditioning is a fact about which couplings survive which updates, not a fact about producers.

\emph{The theorem is nearly analytic.} Nearly, and deliberately: the substance lives in the product form of the state, which the multiplicative update supplies, and in the regime classification, both stated rather than hidden in the proof; the pairwise premise is the one-line case, Theorem~\ref{thm:factorization}, of Theorem~\ref{thm:reichenbach}. The parity witness shows the premise is not vacuous: higher-arity material exists and fails the factorization.

\emph{A faithfulness defender will ask where the independences come from.} From an identity, in the exhibited case: the no-signaling of the two-tag construction is algebraic \cite{AethusQuant}, not a protected cancellation. The defender's question remains legitimate for the model class of this paper: an arrowless nonnegative coupling model does not protect no-signaling by itself. Perturb the singlet factor to $G_\eta = [1 - ab\,\mathbf{x}\cdot\mathbf{y}]\,[1 + \eta a f(\mathbf{y})]$ with $|\eta f(\mathbf{y})| < 1$: summing over $b$ gives $S_\eta = 2[1 + \eta a f(\mathbf{y})]$, and normalization gives $P_\eta(a \mid \mathbf{x}, \mathbf{y}) = \tfrac{1}{2}[1 + \eta a f(\mathbf{y})]$, which signals whenever $f$ varies with the far setting. What protects no-signaling is an independently justified restriction excluding such couplings, imported from the companion quantum construction and its permitted local operations, as the interpretive restraint of Section~\ref{sec:principle} says.

\emph{Interventions distinguish what observation cannot.} They do, and the paper's position on this is stated once rather than staged as a dilemma. $A \to C \to B$, $A \leftarrow C \to B$ and $A \leftarrow C \leftarrow B$ share one screening structure, and interventions distinguish them: settling $A$ moves $B$ only under the first, settling $B$ moves $A$ only under the third, and under the second neither moves the other (Corollary~\ref{cor:inherited}). The framework reproduces those profiles because it accepts the directed structure that distinguishes them: the binding modes are read off the supplied mechanism history, the input sets and normalized mechanisms, which determine a directed acyclic graph, and with them settling is the do-surgery of every causal Bayesian network (Theorem~\ref{thm:settling}). The binding modes are not derived from the joint coupling, and no theory derives interventional laws from the joint coupling alone (Proposition~\ref{prop:parity}); the do-calculus takes its graph as input for the same reason. What the paper declines is only the reading of the directed structure as primitive production: the orientation is the supplied mechanism input structure, carried for its indexing and its counterfactual semantics alike, not an annotation of producers and not an order of acquisition. Whether that interpretation makes production dispensable is the philosophical question of Section~\ref{sec:concept}, and the correspondence theorems do not settle it.

\section{Conclusion} \label{sec:conclusion}

The results are the three groups of the introduction, each with the limitation stated beside it: the observational theorems of Section~\ref{sec:theorem}, which establish factorization consequences and not causal common ancestry; the correspondences of Section~\ref{sec:rungs}, conditional on a supplied mechanism history and, for counterfactuals, on functional entries, which retain directed mechanism structure while declining to interpret it as primitive production; and the boundary of Section~\ref{sec:boundary}, where Bell-violating statistics admit no global noncontextual, measurement-independent valuation, the framework's record discipline explaining, as a premise about its admissible states, why no such mediator is carried. Interpreting these results as a replacement for causal production is a philosophical proposal, stated as one, and its center is the ontological claim that an indexed state describes reality relative to its index completely, with its hard structure not read off its table; Example~\ref{ex:twostates} is the consequence of that claim the paper exhibits, and whether indexing also yields a theorem or constraint that the standard structures do not is a question for further work.

\noindent\textit{Computational checks.} Every enumeration reported in the text is reproduced by the supplementary script \texttt{causality\_checks.py} and its five seeded component scripts (connectivity, separation and screening over 1{,}500 random product-form states; settling against the truncated factorization and the back-door adjustment over 1{,}500 random networks; counterfactual settling against a twin-world computation over 3{,}000 random functional models; far-wing balance and the singlet; the regularity proposition over 600 random deterministic states), together with the three examples added in this revision: the setting-independent screener with CHSH value $2\sqrt2$, the three-valued screener with covariance $\tfrac{4}{75}$, and the two states of Example~\ref{ex:twostates}. The script's expected output accompanies it.

\section*{Provenance}
\addcontentsline{toc}{section}{Provenance}

\noindent AI assistance was used in three ways. Section~\ref{sec:rerendition} was drafted by Claude Fable 5.1 on 2026-09-26 at the author's request and corrected the same day after a referee supplied the counterexamples reproduced in the text, the author's reading, that correlation is weight and that hard invalidity and weight together express the correlational and causal picture, being retained in the form the representation theorem supports; mathematical checking, including the enumerations reported above; and pre-submission audit. Responsibility for every claim is the author's, and independent human refereeing precedes any submission. The analysis extracts, and in places compresses, material developed in the framework's long-form treatment \cite{Aethus}; the behavioral definition of the horns and the factorization display were forced by a parity-style counterexample to a topological formulation, kept in the text as a working part of the result, and the tickle-as-certificate identity of Section~\ref{sec:instances} emerged from the decision-theoretic side of the same development \cite{AethusDecision}. The same primitive enters the decision extraction's sorting of the classical decision-theoretic cases and the quantum--classical classification \cite{AethusDecision, AethusQuant}; that one structure enters all three is offered as a reason to think the proposal is tracking something structural, and not as an argument.

\printbibliography

\end{document}